\uuid{o2WL}
\titre{Aire/volume}
\auteur{bodin ; hanani ; mzari}
\organisation{exo7}
\licence{CC BY-NC-SA 4.0}
\datecreate{2020-01-20}
\niveau{L1}
\module{Géométrie}
\chapitre{Géométrie de l'espace}
\sousChapitre{Aire, volume}
\difficulte{2}
\theme{Aire -- Volume, aire, volume}
\exo7qcmid{lille-1.geometrie-espace.003}

\contenu{
\texte{
Pour ces questions, l'espace est muni d'un repère orthonormé direct  $(O, \vec{i}, \vec{j}, \vec{k})$.
}

\begin{question}
\qclassification{141.02}
\qtags{motcle=aire/volume}
Soit $\vec{u}(-1,1,1), \vec{v}(0,1,2)$ et $\vec{w}(1,0,-1)$ trois vecteurs. Quelles sont les assertions vraies ?
\begin{answers}

      \bad{L'aire du parallélogramme engendré par $\vec{u}$ et  $\vec{v}$ est : $\sqrt 3$.}

       \good{L'aire du parallélogramme engendré par $\vec{u}$ et  $\vec{v}$ est : $ \sqrt 6$.}

    \bad{Le volume du parallélépipède engendré par $\vec{u}$, $\vec{v}$ et  $\vec{w}$ est $1$.}

    \good{Le volume du parallélépipède engendré par $\vec{u}$, $\vec{v}$ et  $\vec{w}$ est $2$.}

\end{answers}
\begin{explanations}
L'aire du parallélogramme engendré par deux vecteurs $\vec{u}$ et $\vec{v}$ est donnée par : $ \Vert \vec{u} \wedge \vec{v} \Vert$.
Le volume du parallélépipède engendré par trois vecteurs $\vec{u}$, $\vec{v}$ et  $\vec{w}$ est donné par : $|\det (\vec{u}, \vec{v},\vec{w})|$.
\end{explanations}
\end{question}
}
